% !TeX root = ../main.tex

\section{Discussion}
\label{sec:discussion}

Bayesian-PMALA brings exact posterior sampling to quantile trend filtering while keeping the proximal structure of the optimization problem. Under the normal--exponential representation, the proximal map of the whole conditional posterior of the trend is a weighted trend filtering problem, so no part of the check loss is smoothed and the Metropolis--Hastings correction is exact. In the simulation study each competitor is most accurate on the benchmark that suits it, and Bayesian-PMALA is competitive with both on both signals, with one default configuration for every signal, noise family and quantile level. It matches the locally adaptive horseshoe sampler with a single global penalty per quantile, and it matches ADMM while learning the penalty and the likelihood scale from the posterior instead of selecting a penalty by an information criterion.

The computational structure is the second difference from the Gibbs sampler of \citet{onizuka2024fast}. Its supplied implementation forms a dense inverse and uses dense Gaussian simulation at each iteration, and our audit found that its sparse option still enters this dense path. Every linear system in Bayesian-PMALA is banded, with bandwidth $k+1$ for every $n$. The metric $M$ and the ADMM system matrix change with $(z, \lambda, \sigma)$ and are therefore factorized at each sweep, but one factorization of each serves the forward and reverse proximal maps and the proposal noise, and for a fixed ADMM iteration budget the work per sweep grows linearly with $n$. Cost per effective posterior draw also depends on mixing, and Section~\ref{sec:scalability} reports measured costs under a common protocol.

The experiment of Section~\ref{sec:air_simulation} clarifies what the global penalty buys. There the estimand is a peak-free baseline at a low quantile and the observed series is contaminated by transient peaks, so local adaptivity is a liability: the horseshoe's coefficient-specific scales let the fit follow each peak, which lowers the check loss on the contaminated data, raises the error against the peak-free target, and removes from the detrended residuals part of the signal the downstream detector is meant to find. A single penalty cannot bend locally, and Bayesian-PMALA is far more accurate than the horseshoe sampler on that design and close to ADMM. The design also shows the limit of what the likelihood can identify: the separation between a smooth background and transient bursts is not in the data model $y = \theta + \varepsilon$, so it has to enter through the prior on $\lambda$. We use a second prior, fixed in advance and shared across sample sizes, for data of this kind, and keep the weak default elsewhere; a hierarchy that learns the separation itself, for instance by modelling the burst component explicitly, is a natural extension.

Treating the ALD scale $\sigma$ as unknown is a substantive modeling choice. At fixed $(\lambda, \sigma)$ the conditional mode is the quantile trend filter with penalty $\lambda\sigma$ \eqref{eq:map_qtf}, so the posterior learns the effective penalty through two parameters, the scale adapting to the noise level and $\lambda$ to the roughness of the trend, rather than asking one tuning parameter to absorb both roles. Learning $\sigma$ does not by itself guarantee frequentist coverage when the asymmetric Laplace likelihood is misspecified, a point also recognized by the calibration approach of \citet{onizuka2024fast}.

A single global penalty is a deliberate simplification relative to coefficient-specific shrinkage, and the experiments do not establish a universally preferred prior or difference order. The pointwise intervals are below nominal coverage at the outer quantile levels (Section~\ref{sec:sim:ind:smooth}), and their calibration there requires further study.

% TODO (2026-09-26): the next paragraph concerns the joint sampler of the previous version; rewrite with Section~\ref{subsec:joint}.
The sentinel quantile construction used in Section~\ref{sec:sim:extreme} --- auxiliary quantile levels placed beyond the target range and sampled jointly so that the box constraints of Section~\ref{subsec:joint} bound each extreme target dynamically --- resolves the gradient degeneracy of P-MALA at extreme quantile levels and produces stable estimates in the simulation study. It is, however, a practical device rather than a formally justified procedure. The specific recommendations --- up to five sentinels per side, equally spaced within each side, with spacings chosen to fill the available probability room --- are motivated by empirical experience and by the requirement that the sentinels be close enough to the target to provide a tight dynamic bound while remaining far enough apart to avoid degeneracy among themselves.

Several extensions follow from the present work. Non-crossing constraints enter the proximal map as box constraints, so a joint version for several quantile levels keeps the banded structure (Section~\ref{subsec:joint}). Locally adaptive priors can be combined with the same proximal update, since the conditional of the trend remains a weighted quadratic plus a weighted $\ell_1$ penalty. Extension to spatially indexed signals requires generalizing the difference operator to graphs, for which the systems are no longer banded and sparse or multigrid solvers may serve a similar role. Finally, the convergence properties of the full Gibbs sampler with continuous step-size adaptation remain to be established.
